\section{Boost Converter Modeling}
\label{sec:converter-model}

\subsection{Specifications and Component Sizing}
The MPPT controllers command the duty cycle of a boost converter interposed
between the PV module/string and a resistive load representing a downstream
DC bus, switching at $f_s=20\,$kHz with an input range
$V_{in}\in[20,40]\,$V (covering the module's expected $V_{mp}$ range under
varying conditions) and output $V_{out}=48\,$V. Sizing follows the standard
CCM boost relations
\begin{align}
D &= 1 - \frac{V_{in}}{V_{out}}, \\
L &= \frac{V_{in} D}{f_s \, \Delta I_L}, \qquad \Delta I_L = 0.30\, I_{L,avg}, \\
C &= \frac{I_{out} D}{f_s \, \Delta V_o}, \qquad \Delta V_o = 0.01\, V_{out},
\end{align}
with the 30\% inductor-current-ripple and 1\% output-voltage-ripple design
targets following \cite{erickson2001fundamentals}. At the worst-case
operating point ($V_{in}=20\,$V, full load $P_{max}=250\,$W), this gives
$D=0.583$, $L=155.6\,\mu$H, $C=316.5\,\mu$F, and load resistance
$R=V_{out}^2/P_{max}=9.22\,\Omega$.

\subsection{Small-Signal Model}
The state-space averaged control-to-output transfer function for a CCM
boost converter is
\begin{equation}
G_{vd}(s) = \frac{V_{out}}{(1-D)^2} \cdot
\frac{1 - s/\omega_z}{1 + \dfrac{s}{Q\omega_0} + \left(\dfrac{s}{\omega_0}\right)^2},
\label{eq:small-signal}
\end{equation}
with
\begin{equation}
\omega_0 = \frac{1-D}{\sqrt{LC}}, \quad
Q = (1-D)R\sqrt{\frac{C}{L}}, \quad
\omega_z = \frac{(1-D)^2 R}{L}.
\end{equation}
At the design operating point, $\omega_0 = 1878\,$rad/s (299\,Hz),
$Q=5.48$, and the right-half-plane (RHP) zero $\omega_z=10{,}286\,$rad/s
(1637\,Hz). The RHP zero is included explicitly in \eqref{eq:small-signal}
because a boost converter's control-to-output response has non-minimum-phase
behavior (an initial dip in output before the commanded rise) that a
simplified transfer function omitting the zero would misrepresent -- of
direct relevance to any MPPT algorithm that assumes a monotonic response to
a duty-cycle perturbation.

\subsection{Discontinuous Conduction Mode (DCM)}
The critical inductance for continuous conduction at a given operating point
is $L_{crit} = (1-D)^2 D R / (2 f_s)$. At full load this design's
$L=155.6\,\mu$H exceeds $L_{crit}=23.3\,\mu$H, confirming CCM operation; at
light load (5\% of $P_{max}$, representative of low-irradiance conditions),
the same inductor value no longer satisfies the CCM condition and the
converter enters DCM.

This has a concrete implication for MPPT behavior that is not simply a
converter-design footnote: in DCM, the $V_{out}/V_{in}$ relationship depends
on $L$, $f_s$, and load in addition to duty cycle, so \eqref{eq:small-signal}
no longer holds and the same duty-cycle perturbation produces a different,
load-dependent voltage/power response than in CCM. Every algorithm in
Section~\ref{sec:algorithms} infers tracking direction from a duty-cycle
perturbation's effect on $dP/dV$ (P\&O, IncCond, fuzzy logic, and SMC's
sliding surface directly; Q-learning's reward indirectly) without
re-deriving or re-tuning against a DCM model. A spurious sensitivity change
at light load is therefore a plausible confound for any algorithm whose
convergence or oscillation metrics degrade specifically in this paper's
low-irradiance scenarios (200\,W/m$^2$ in the multi-level-irradiance and
rapid-double-step scenarios, or the low-irradiance modules in the
partial-shading scenarios) -- a possibility we return to in
Section~\ref{sec:conclusion}'s discussion of limitations, not something
addressable by changing the MPPT algorithm itself.
